% ===============================================================
% GRADE 9 -- Real-world multiple-choice problems
% ===============================================================
\subsection{Real-World Problems (Grade 9)}
\probset{9}
{\small Read the situation, pick a letter within 20 seconds, then check the Answer Key.\par}
\begin{gradebody}

\textbf{Term 1}

\mcq{A tricycle driver ties sacks of rice tightly to the roof. If they were loose, they would fly forward when he brakes suddenly. Why?}
{Gravity pulls them forward}{The brakes push them}{Air resistance}{Inertia keeps them moving}
{D}{Objects keep moving unless a force (here, the rope) stops them.}{A}

\mcq{The same push is given to an empty grocery cart and a full one. The empty cart speeds up more because it has}
{more inertia}{more friction}{less mass}{more momentum}
{C}{For the same force, $a = F/m$ is larger when mass is smaller.}{E}

\mcq{A \SI{1000}{kg} car accelerates at \SI{2}{\meter\per\second\squared}. What net force acts on it?}
{\SI{2000}{N}}{\SI{998}{N}}{\SI{500}{N}}{\SI{1002}{N}}
{A}{$F = ma = 1000 \times 2$.}{A}

\mcq{A fisherman jumps from his bangka onto the pier, and the bangka drifts backward. Which law explains this?}
{Third law}{First law}{Law of gravitation}{Second law}
{A}{He pushes the boat back; the boat pushes him forward.}{E}

\mcq{You push on a concrete wall with \SI{50}{N}. The wall pushes on you with}
{more than \SI{50}{N}}{\SI{50}{N}}{less than \SI{50}{N}}{\SI{0}{N}}
{B}{Action and reaction are equal and opposite.}{E}

\mcq{A rocket in space, with no air to push against, can still speed up because}
{space has wind}{gravity pushes it}{the exhaust gas pushes back on the rocket}{inertia increases}
{C}{The rocket pushes gas backward; the gas pushes the rocket forward (third law).}{D}

\mcq{At home, turning off the TV does not turn off the electric fan. The outlets are wired in}
{a single loop}{series}{parallel}{a short circuit}
{C}{Parallel branches work independently.}{E}

\mcq{A flashlight with two \SI{1.5}{V} cells in series is brighter than one with a single cell because}
{the cells share the bulb}{the voltages add to \SI{3}{V}}{the resistance drops to zero}{the current path is longer}
{B}{Series cells add their voltages, pushing more current through the bulb.}{A}

\mcq{A rice cooker draws \SI{4}{A} from a \SI{220}{V} outlet. Its resistance is}
{\SI{224}{\ohm}}{\SI{880}{\ohm}}{\SI{55}{\ohm}}{\SI{0.018}{\ohm}}
{C}{$R = V/I = 220/4$.}{A}

\mcq{Plugging an aircon, flat iron, and kettle into one extension cord trips the breaker. Why?}
{Voltage becomes too high}{Resistance becomes infinite}{The cord becomes a series circuit}{Total current becomes too high (overload)}
{D}{Each parallel load adds current through the cord.}{A}

\mcq{Why is touching an outlet with wet hands especially dangerous?}
{Water raises the voltage}{Water lowers skin resistance, so more current flows}{Water is an insulator}{Water stores charge}
{B}{By $I = V/R$, lower resistance means a bigger current.}{D}

\mcq{A bird sits on a single high-voltage wire without being shocked because}
{its feet are insulated}{there is almost no voltage difference across its body}{birds have high resistance}{the wire carries no current}
{B}{Current needs a potential difference; both feet are at the same voltage.}{D}

\mcq{Airport baggage scanners see inside bags using}
{X-rays}{ultraviolet}{radio waves}{infrared}
{A}{X-rays pass through soft materials but not dense objects.}{E}

\mcq{After an earthquake, rescuers use a thermal camera to find trapped survivors under rubble. It detects}
{infrared}{gamma rays}{ultraviolet}{microwaves}
{A}{Warm bodies give off infrared radiation.}{A}

\mcq{Students get sunburned at a beach outing even on a cloudy day because}
{UV passes through clouds}{clouds give off UV}{the sea makes X-rays}{sand gives off infrared}
{A}{Much UV gets through thin clouds.}{A}

\mcq{A microwave oven heats leftover adobo mainly by}
{making water molecules vibrate}{UV radiation}{conduction from the plate}{giving off X-rays}
{A}{Microwaves make water molecules move faster (heat).}{D}

\mcq{FM radio (around \SI{100}{MHz}) and AM radio (around \SI{1}{MHz}): which has the longer wavelength?}
{Equal}{FM}{AM}{Cannot tell}
{C}{Lower frequency means a longer wavelength ($v = f\lambda$).}{D}

\mcq{Why are repeated X-rays limited, while radio waves around us are considered safe?}
{Radio waves are invisible}{X-rays are slower}{X-rays are ionizing and can damage DNA}{X-rays have longer wavelengths}
{C}{High-energy radiation can break chemical bonds in cells.}{A}

\mcq{Fossils of the same freshwater reptile are found in Brazil and South Africa. This suggests that}
{the fossils were planted}{the reptile swam the Atlantic}{the climates are identical}{the continents were once joined}
{D}{This is fossil evidence for continental drift.}{E}

\mcq{Iceland is slowly splitting, and new crust forms in its rift valley. Iceland lies on a}
{hotspot only}{divergent boundary}{convergent boundary}{transform boundary}
{B}{It sits on the Mid-Atlantic Ridge, where plates pull apart.}{A}

\mcq{The deep Philippine Trench east of Mindanao formed because}
{a river eroded it}{plates slide past each other}{plates pull apart}{an oceanic plate subducts}
{D}{Trenches mark convergent subduction zones.}{A}

\textbf{Term 2}

\mcq{In a road cut, trilobite fossils lie in a layer \emph{below} a layer with dinosaur fossils. Which is older?}
{Cannot tell}{The trilobite layer}{The dinosaur layer}{They are the same age}
{B}{By superposition, lower layers are older.}{A}

\mcq{The best way to date a wooden tool from an old burial site is}
{carbon-14 dating}{counting rock layers}{index fossils}{uranium--lead dating}
{A}{C-14 dates once-living material up to tens of thousands of years old.}{A}

\mcq{After a big earthquake, stations on the opposite side of Earth record P-waves but no S-waves. What does this tell us?}
{The mantle is liquid}{The inner core is hollow}{The outer core is liquid}{The crust is thin}
{C}{S-waves cannot travel through liquid.}{D}

\mcq{A compass needle points north because of Earth's magnetic field, which is produced by}
{the Moon}{the atmosphere}{the crust}{moving molten iron in the outer core}
{D}{Convection of liquid iron in the outer core generates the field.}{A}

\mcq{During the Geminids, you see many bright streaks in the night sky. They are caused by}
{falling stars}{meteoroids burning up in the atmosphere}{comets hitting Earth}{satellites}
{B}{A meteor is the glow of a burning meteoroid.}{A}

\mcq{Villagers find a hot, dark rock in a rice field after seeing a fireball. It is a}
{asteroid belt}{meteor}{comet}{meteorite}
{D}{A meteoroid that reaches the ground is a meteorite.}{E}

\mcq{A comet's tail always points away from the Sun because}
{solar wind and radiation push the gas and dust}{of its orbit direction}{the comet is spinning}{of gravity from Jupiter}
{A}{The tail follows the Sun's outward push, not the direction of travel.}{D}

\mcq{The Hubble telescope orbits above the atmosphere because}
{it is closer to the stars}{it is easier to repair}{the air blurs and absorbs light}{it needs sunlight for power}
{C}{Above the air there is no blurring and no blocked wavelengths.}{A}

\mcq{The Philippine microsatellite Diwata-1 was mainly used for}
{a Moon landing}{communication with Mars}{human spaceflight}{Earth observation (weather, crops, disasters)}
{D}{It imaged the Earth to help agriculture and disaster response.}{A}

\mcq{A forest with many tree species recovers from a pest outbreak faster than a one-species plantation because}
{forests have no animals}{pests prefer forests}{high biodiversity gives alternatives and resilience}{plantations grow faster}
{C}{The pest can't wipe out every species at once.}{A}

\mcq{Janitor fish released into Laguna de Bay harm native fish mainly because they}
{eat people}{make the water salty}{are poisonous}{are invasive and compete with few natural predators}
{D}{Invasive species spread unchecked and disrupt the ecosystem.}{A}

\mcq{The first effect of kaingin on a mountain forest is}
{more biodiversity}{habitat loss and soil erosion}{more rainfall}{cooler temperatures}
{B}{Clearing removes habitat and the roots that hold the soil.}{E}

\mcq{A coastal barangay replants bakawan along its shoreline to reduce damage from storm surges. They are planting}
{seagrass only}{mangroves}{coconut trees}{coral reefs}
{B}{Mangrove roots weaken waves and hold the sediment.}{E}

\mcq{Corals in Tubbataha turn white during an unusually warm year. This is}
{coral bleaching from heat stress}{sand covering}{normal growth}{coral spawning}
{A}{Heat makes corals expel their algae (zooxanthellae).}{A}

\mcq{Philippine eagles are bred at a centre in Davao. This is an example of}
{in-situ conservation}{extinction}{habitat restoration}{ex-situ conservation}
{D}{Conservation outside the natural habitat is ex-situ.}{A}

\textbf{Term 3}

\mcq{Police compare DNA from a strand of hair with a suspect's DNA. This works because}
{each person's DNA is unique (except identical twins)}{all humans have identical DNA}{hair has no DNA}{DNA changes daily}
{A}{DNA sequences differ between individuals.}{A}

\mcq{Smoking and heavy sun exposure raise cancer risk because they}
{cause mutations in DNA}{increase cell size}{add vitamins}{stop mitosis}
{A}{Mutagens (chemicals, UV) can damage DNA.}{E}

\mcq{A DNA strand reads ATG CCA. Its complementary strand reads}
{GCA TTG}{ATG CCA}{TAC GGT}{UAC GGU}
{C}{A pairs with T and G pairs with C (no U in DNA).}{A}

\mcq{A mutation lets a bacterium survive an antibiotic. For the bacterium, this mutation is}
{beneficial}{neutral}{impossible}{harmful}
{A}{It improves survival in that environment.}{A}

\mcq{Salt water conducts electricity, but sugar water does not. Salt is}
{covalent}{ionic}{metallic}{a mixture}
{B}{Dissolved ionic compounds release ions that carry charge.}{A}

\mcq{Copper is used for house wiring because metallic bonding makes it}
{transparent}{brittle}{an insulator}{a good conductor and ductile}
{D}{Free-moving electrons conduct, and the metal can be drawn into wires.}{E}

\mcq{Breathing out through a straw into clear limewater makes it milky. This shows exhaled air contains}
{oxygen}{carbon dioxide}{water vapour}{nitrogen}
{B}{\ce{CO2} forms a \ce{CaCO3} precipitate with limewater.}{E}

\mcq{Magnesium loses 2 electrons and chlorine gains 1. In magnesium chloride, the ratio Mg : Cl is}
{2:1}{1:1}{2:3}{1:2}
{D}{One \ce{Mg^{2+}} balances two \ce{Cl-}, giving \ce{MgCl2}.}{D}

\mcq{Table salt melts at about \SI{800}{\celsius}, but candle wax melts below \SI{100}{\celsius}. Why?}
{Salt is a metal}{Wax has no bonds}{Salt's ionic lattice has strong attractions throughout}{Salt has fewer atoms}
{C}{Ionic lattices need much more energy to break than weak forces between molecules.}{D}

\end{gradebody}
\probstats
